\section{Related Work}
\label{sec:related}

\subsection{Long- and Ultra-Long Video Understanding}
\label{sec:related-long-video}

Long-video understanding requires balancing temporal coverage against the cost of retaining detailed visual evidence. One line of work expands the video context available to multimodal models through compact representations. MovieChat progressively consolidates dense video tokens into sparse memory~\cite{song2024moviechat}, while LongVU reduces temporal and spatial redundancy through adaptive compression~\cite{shen2025longvu}. Video-XL summarizes visual information into compact key-value representations for hour-scale inputs~\cite{shu2025videoxl}, and VideoChat-Flash compresses video tokens hierarchically at the clip and video levels~\cite{li2026videochatflash}. These methods balance compression against the retention of question-relevant details. At week-long temporal horizons, EgoMemReason evaluates reasoning over egocentric recordings and emphasizes the retrieval and integration of temporally distant evidence~\cite{wang2026egomemreason}.

Another line of work externalizes video content into searchable memories and retrieves evidence on demand. DrVideo represents videos as text documents and iteratively retrieves and augments relevant evidence~\cite{ma2025drvideo}. M3-Agent organizes audiovisual observations into entity-centric episodic and semantic memories~\cite{long2026m3agent}, while WorldMM adaptively retrieves from episodic, semantic, and visual memories at multiple temporal granularities~\cite{yeo2026worldmm}. More recently, MERIT builds a multi-key episodic index and expands the temporal context around retrieved segments at inference time~\cite{choi2026merit}. These methods shift part of the challenge from compressing visual context to selecting and combining evidence from external memory.

Complementary to memory design, agentic methods focus on iterative tool use for evidence acquisition. Ego-R1 uses supervised fine-tuning and reinforcement learning to train a controller that invokes hierarchical retrieval and visual perception tools for ultra-long egocentric video reasoning~\cite{tian2026egor1}. VideoPro adaptively chooses between direct video-language reasoning and executable programs, refining program-based reasoning with execution feedback and confidence signals~\cite{li2026videopro}. Structured querying also has direct precedents: DoraemonGPT builds symbolic video memories queried through SQL-generating tools~\cite{yang2024doraemongpt}, while VideoAgent combines temporal memory with an object-memory database accessed through agent-generated SQL~\cite{fan2025videoagent}.

VideoSQL builds on these approaches by coupling SQL-based access to video-derived indices with trainable skills. These skills capture reusable procedures for querying and combining indexed evidence across questions in ultra-long video understanding.
